\documentclass[11pt,a4paper]{article}
\usepackage[utf8]{inputenc}
\usepackage[T1]{fontenc}
\usepackage{amsmath, amssymb, amsthm}
\usepackage{geometry}
\geometry{margin=2.5cm}
\usepackage{lmodern}
\usepackage[final,expansion=false]{microtype}
\usepackage{booktabs}
\usepackage{array}
\usepackage{enumitem}
\usepackage{xcolor}
\usepackage{titlesec}
\PassOptionsToPackage{hyphens}{url}
\usepackage{hyperref}
\usepackage{xurl}
\hypersetup{
  colorlinks=true,
  linkcolor=blue!50!black,
  urlcolor=blue!50!black,
  citecolor=blue!50!black,
  pdftitle={The Universal Right-Triangle Reflection Sequence},
  pdfauthor={Vico Bonfioli},
}
\usepackage{parskip}
\setlength{\parindent}{0pt}
\setlength{\parskip}{0.55em}

\titleformat{\section}{\normalfont\Large\bfseries\color{blue!40!black}}{\thesection.}{0.6em}{}

\newtheorem{theorem}{Theorem}
\newtheorem{lemma}[theorem]{Lemma}
\newtheorem{corollary}[theorem]{Corollary}
\newtheorem{proposition}[theorem]{Proposition}
\newtheorem{example}[theorem]{Example}
\theoremstyle{remark}
\newtheorem{remark}[theorem]{Remark}
\theoremstyle{plain}

\title{\textbf{The Universal Right-Triangle Reflection Sequence (Unequal Legs)}\\[0.3em]
\large A closed-form rate $r_n = 1 + 2\cos(2\pi/(n+3))$ and unconditional Class~C faithfulness in every dimension}
\author{Vico Bonfioli\thanks{Independent researcher, Italy.
\texttt{vicobonfioli@gmail.com}.\,
\url{https://elvec1o.github.io/home/posts/}.}}
\date{\today}

\begin{document}
\maketitle

\begin{abstract}
For every right triangle $T$ in the plane with positive \textbf{unequal} rational legs, the BFS orbit-growth sequence $u_d^T$ under the group of side-reflections is independent of $T$ (Theorem~\ref{thm:universality}); we call the common sequence the \emph{universal right-triangle reflection sequence}.  The kernel of the side-reflection representation is identified as a cyclic left ideal $\mathbb{Z}[Q]\cdot h$ on the explicit central generator $h = (c-1)(Y-1)$ via Mayer--Vietoris and Crowell.  In every dimension $n \ge 2$, the right-corner orthoscheme reflection group has closed-form Coxeter--Steinberg growth rate $r_n = 1 + 2\cos(2\pi/(n+3))$ (Theorem~\ref{thm:ndim-rate}).  For pairwise distinct legs the affine representation is unconditionally faithful in every dimension $n \ge 3$ (Theorem~\ref{thm:master-C-uncond}), via a uniform Schur-complement determinant identity $\det Q_n = -\prod_i a_i^2$ (Lemma~\ref{lem:uniform-detQ}) together with Mordell descent on the two Cremona curves $\mathtt{24a1}$ and $\mathtt{72a2}$.  A dimension-independent collision-depth invariant $\mathrm{cd}_n(\lambda) \in \{\infty, 4, 3\}$ determined by the $(e, t)$-partition statistics of the leg sequence (Theorem~\ref{thm:cd-general}) classifies first-collision behaviour across all dimensions.  The eight length-$10$ relations, the Schur-complement identity, and the universality of the orbit-growth sequence through depth $22$ are formally machine-verified in the Lean~4 proof assistant (\S\ref{sec:reproduce}).  Conditional partial-progress results on the asymptotic ratio $\beta_2 \approx 1.47$ of the 2D universal sequence, together with the mod-$p$ state-space and Knuth--Bendix experiments, are deferred to the companion document \texttt{paper\_extra.tex}.
\end{abstract}

\section{Setting}
\label{sec:setting}

Let $T$ be a triangle with vertices $V_0, V_1, V_2 \in \mathbb{Q}^2$.
Only the vertex coordinates and the leg lengths $a, b$ are assumed
rational.  For $i \in \{0, 1, 2\}$ write $R_i$ for the reflection of
$\mathbb{R}^2$ across the line containing the side opposite $V_i$.
Set $W(T) = \langle R_0, R_1, R_2 \rangle \le \mathrm{Isom}(\mathbb{R}^2)$.

When we need explicit coordinates (from §\ref{sec:bonfioli} onward) we
place the right-angle vertex at $B = (a, 0)$ with $A = (0, 0)$,
$C = (a, b)$, and write $a, b$ for the two leg lengths.

We explore the orbit of $T$ under $W(T)$ by breadth-first search in the
Cayley graph relative to the generating set $\{R_0, R_1, R_2\}$.  Let
\[
\mathcal{O}_d(T) = \bigl\{\, g \cdot T : g \in W(T) \text{ has minimal
word-length } d \,\bigr\},
\]
and set $u_d := |\mathcal{O}_d(T)|$, $u_0 = 1$.
All distances are computed in exact rational arithmetic and the proofs
are checked symbolically over $\mathbb{Q}(a, b)$.

\subsection{Notation}
\label{ssec:notation}

\begin{itemize}[leftmargin=*,itemsep=0.1em]
\item $u_d$: BFS layer size at minimal word-length $d$.
\item $\delta_d := F_{d+3} - u_d$: Fibonacci deficit against the Steinberg/Fibonacci baseline.
\item $\rho_T : W \to \mathrm{Aff}(\mathbb{R}^2)$: the affine action of
$W = \langle R_0, R_1, R_2 \rangle$ induced by the sides of $T$.
\item $\mathrm{cd}_n(\lambda)$: the collision-depth invariant in
dimension $n$ for a leg sequence $\lambda$ (Theorem~\ref{thm:cd-general}).
\item $(e, t)$: positional invariant of a leg sequence; $e$ counts
endpoint-adjacent equal pairs and $t$ counts triples of consecutive
equal legs.
\end{itemize}

\paragraph{Related work.}
The $n$-dimensional rate $r_n = 1 + 2\cos(2\pi/(n+3))$ derived in
\S\ref{sec:ortho} is the closed form expected for the growth rate of a
path-shaped Coxeter group with one $\infty$ edge, in the framework of
Cannon--Floyd--Parry (Cannon--Wagreich; Floyd--Plotnick).  The same
family of algebraic integers also appears in birational surface
dynamics (Diller--Favre; McMullen; Bedford--Kim); the overlap is at the
level of expressions, not of number-theoretic class.

\section{The Coxeter-group identity (Fibonacci phase)}
\label{sec:fib}

Assume $T$ is a non-isosceles right triangle with the right angle at
$V_2$.  Then $R_0$ and $R_1$ are reflections across the two sides
incident to $V_2$, meeting at $\pi/2$, hence $R_0 R_1 R_0 R_1 = \mathrm{id}$.
The remaining two angles are irrational multiples of $\pi$ by Niven's
theorem (using the non-isosceles assumption), so no relation $(R_i R_j)^m = 1$
holds for any other pair.  The abstract group is
\[
W = \bigl\langle\, R_0, R_1, R_2 \,\bigm|\, R_i^2 = 1\,(i = 0,1,2),
\,(R_0 R_1)^2 = 1\,\bigr\rangle.
\]
By Steinberg's word-growth formula the Poincar\'e series is
$W(t) = (1+t)^2/(1 - t - t^2)$, so $[t^d] W(t) = F_{d+3}$ for $d \ge 1$.

\begin{corollary}
For $1 \le d \le 9$ (i.e.\ until the first Euclidean lattice collision at
$d = 10$), $u_d = F_{d+3}$.  (At $d = 0$, $u_0 = 1 \ne F_3 = 2$.)
\end{corollary}

\section{Eight length-$10$ affine relations}
\label{sec:bonfioli}

Exact-rational breadth-first search yields $u_{10} = 225$, a deficit of
$8$ from $F_{13} = 233$.  These $8$ length-$10$ affine relations are
listed in Table~\ref{tab:bonfioli}.

\begin{table}[ht!]
\centering
\small
\renewcommand{\arraystretch}{1.2}
\begin{tabular}{r l l}
\toprule
\# & word A & word B \\
\midrule
1 & $R_0 R_1 R_2 R_0 R_2 R_0 R_1 R_2 R_1 R_2$ &
    $R_2 R_1 R_2 R_0 R_1 R_2 R_0 R_2 R_0 R_1$ \\
2 & $R_0 R_1 R_2 R_1 R_2 R_0 R_1 R_2 R_0 R_2$ &
    $R_2 R_0 R_2 R_0 R_1 R_2 R_1 R_2 R_0 R_1$ \\
3 & $R_0 R_2 R_0 R_1 R_2 R_1 R_2 R_0 R_1 R_2$ &
    $R_2 R_0 R_1 R_2 R_1 R_2 R_0 R_1 R_2 R_0$ \\
4 & $R_0 R_2 R_0 R_2 R_0 R_1 R_2 R_1 R_2 R_0$ &
    $R_1 R_2 R_1 R_2 R_0 R_1 R_2 R_0 R_2 R_1$ \\
5 & $R_0 R_2 R_0 R_2 R_0 R_1 R_2 R_1 R_2 R_1$ &
    $R_1 R_2 R_1 R_2 R_0 R_1 R_2 R_0 R_2 R_0$ \\
6 & $R_0 R_2 R_1 R_2 R_0 R_1 R_2 R_0 R_2 R_0$ &
    $R_1 R_2 R_0 R_2 R_0 R_1 R_2 R_1 R_2 R_1$ \\
7 & $R_0 R_2 R_1 R_2 R_0 R_1 R_2 R_0 R_2 R_1$ &
    $R_1 R_2 R_0 R_2 R_0 R_1 R_2 R_1 R_2 R_0$ \\
8 & $R_1 R_2 R_0 R_1 R_2 R_0 R_2 R_0 R_1 R_2$ &
    $R_2 R_0 R_1 R_2 R_0 R_2 R_0 R_1 R_2 R_1$ \\
\bottomrule
\end{tabular}
\caption{The eight length-$10$ affine relations.  Each row's two words
produce the same affine matrix in the side-reflection action of $W$ on
any right triangle.}
\label{tab:bonfioli}
\end{table}

\begin{theorem}[Eight length-$10$ affine relations]
\label{thm:symbolic}
Each of the eight identities of Table~\ref{tab:bonfioli} holds for
every right triangle.  With the placement $A = (0, 0)$, $B = (a, 0)$,
$C = (a, b)$, $a, b > 0$, the affine matrices $M_A = R_{w_A[0]} \cdots R_{w_A[9]}$
and $M_B = R_{w_B[0]} \cdots R_{w_B[9]}$ satisfy $M_A = M_B$ as
$3 \times 3$ affine matrices over $\mathbb{Q}(a, b)$.
\end{theorem}

\begin{proof}
This is a finite algebraic verification.  Each reflection $R_i$ acts as
an affine isometry whose matrix entries are rational functions of
$a, b$ over $\mathbb{Q}$ (with denominators dividing a power of
$a^2 + b^2$).  Each entry of $M_A - M_B$ is a rational function in
$\mathbb{Q}(a, b)$; clearing denominators gives a polynomial identity
which is verified by SymPy in $\mathbb{Q}(a, b)$ (deterministic, no
heuristics).  All eight identities have additionally been
\emph{formally verified} in the Lean~4 proof assistant: the theorems
\texttt{rel1\_symbolic}, \dots, \texttt{rel8\_symbolic} in
\texttt{lean/with\_mathlib/SymbolicVerification.lean} each discharge the
six affine-matrix-entry identities (over $\mathbb{Q}[a,b]$) defining
$M_A = M_B$ by the \texttt{ring} tactic.  As these proofs use only
\texttt{ring}, they are checked by the Lean kernel, so the identities
hold for every right triangle, not merely on a sample of
specializations.
\qed
\end{proof}

\begin{proposition}[Completeness at depth $10$]
\label{prop:bonfioli-completeness}
Among the $233$ canonical Coxeter words of length $10$ in $W$, exactly
$225$ have pairwise distinct images under $\rho_T$ for every right
triangle $T$ --- $217$ singleton classes and $8$ collision pairs, each
pair being one of the rows of Table~\ref{tab:bonfioli}.  No further
length-$10$ affine identities exist beyond the eight relations and the
abstract Coxeter relations.
\end{proposition}

\begin{proof}
Finite enumeration.  All $233$ canonical Coxeter words of length $10$
are enumerated; the affine matrix of each is computed over $\mathbb{Q}$
on the reference $(3, 4, 5)$ triangle in exact rational arithmetic, the
words are grouped by matrix, and the resulting partition has exactly
$225$ blocks ($8$ of size $2$, $217$ of size $1$).  The eight size-$2$
blocks are cross-checked against six additional triangles
$(5, 12), (8, 15), (7, 24), (1, 2), (2, 3), (1, 7)$, and each block
produces the same affine matrix in all seven cases.  By
Theorem~\ref{thm:symbolic} the eight identities hold symbolically over
$\mathbb{Q}(a, b)$ and are therefore the only length-$10$ collisions.
\qed
\end{proof}

\section{Universality across right triangles}
\label{sec:univ}

By Theorem~\ref{thm:symbolic} the kernel $N$ of $W \to \mathrm{Aff}(\mathbb{R}^2)$
contains the eight length-$10$ relations in addition to the Coxeter
relations.  The quotient $W/N$ --- and hence its orbit-growth sequence
$u_d$ --- is therefore independent of the choice of right triangle.

\begin{theorem}[Universality]
\label{thm:universality}
Let $W = \langle R_0, R_1, R_2 \mid R_i^2 = 1,\, (R_0 R_1)^2 = 1\rangle$,
and for each right-triangle parameter $T = (a, b)$ with positive rational legs satisfying
\[
\boxed{\,a \ne b\,}
\]
let $\rho_T : W \to \mathrm{Aff}(\mathbb{R}^2)$ be the side-reflection
realization.  Then $\ker(\rho_T)$ is the same subgroup of $W$ for every
such $T$.  In particular the orbit-growth sequence
$u_d^T := |\{\rho_T(w) : |w|_W = d\}|$ is independent of $T$.
\end{theorem}

\begin{remark}[The isosceles case is genuinely different]
\label{rem:isosceles-excluded}
The hypothesis $a \ne b$ is essential.  When $a = b$, the leg-swap symmetry
adds a $\mathbb{Z}/2$ stabilizer to the orbit action, and the orbit sequence
becomes $1, 3, 5, 8, 11, 13, 16, \ldots$, which is strictly smaller than
the universal sequence $a(n)$ of Table~\ref{tab:tail} for $n \ge 4$.  The
universality phenomenon of Theorem~\ref{thm:universality} does not apply
to the isosceles case.  See the OEIS comment on \texttt{A\_2D\_universal}.
\end{remark}

\begin{remark}[Formal verification of universality through depth $22$]
\label{rem:lean-universality}
The structural argument below proves Theorem~\ref{thm:universality} for
all depths.  Independently, the universality of the orbit-growth sequence
has been \emph{machine-verified} for depths $0 \le d \le 22$ in Lean~4
(theorem \texttt{universal\_layers\_through\_22} in
\texttt{lean/with\_mathlib/ComputableUniversality.lean}).  That proof
carries out a single breadth-first sweep over the canonical Coxeter words
of $W$, deduplicating the symbolic affine images
$\rho_{(a,b)}(w) \in \mathrm{Aff}(\mathbb{Q}(a,b))$ by their numerators
relative to the common denominator $(a^2+b^2)^{22}$, and certifies that
the layer counts are exactly
\[
1, 3, 5, 8, 13, 21, 34, 55, 89, 144, 225, 351, 554, 875, 1345,
2066, 3203, 4971, 7574, 11543, 17683, 27108, 41067
\]
\emph{simultaneously for every right triangle with positive unequal
legs}.  This kernel-checked certificate matches the depth reached by the
independent multi-triangle reproduction of \S\ref{sec:tail} and removes
any reliance on a single specialization for $d \le 22$.
\end{remark}

The argument factors through three lemmas.

\begin{lemma}[Mayer--Vietoris reduction]
\label{lem:trans-mod}
Let $A = V_4 = \langle R_0, R_1 \mid R_0^2, R_1^2, (R_0 R_1)^2\rangle$,
$B = \mathbb{Z}/2 = \langle R_2 \mid R_2^2\rangle$, so $W = A * B$.  Let
$Q = D_\infty \times \mathbb{Z}/2$, with generators $X, X', Y, c$
subject to $X^2 = (X')^2 = Y^2 = c^2 = 1$, $X X' = c$, and $c$ central.
Let $T_\rho = \ker(W \twoheadrightarrow Q)$.  Then, as left
$\mathbb{Z}[Q]$-modules,
\[
T_\rho^{\mathrm{ab}} \;\cong\; H_1(W;\, \mathbb{Z}[Q])
   \;\cong\; I_A \cap I_B \;=\; \mathbb{Z}[Q] \cdot h,
\]
where $h := (c-1)(Y-1) \in \mathbb{Z}[Q]$.  The annihilator of $h$ is
generated by $(Y+1)$ and $(c+1)$:
$(Y+1)\, h = 0$ and $(c+1)\, h = 0$ in $\mathbb{Z}[Q]$, and
\[
T_\rho^{\mathrm{ab}} \;\cong\; \mathbb{Z}[Q]/(Y+1,\, c+1) \;=:\; M.
\]
\end{lemma}

\begin{proof}
Since $A \hookrightarrow Q$ and $B \hookrightarrow Q$ are injective,
$\mathbb{Z}[Q]$ is free as $\mathbb{Z}[A]$- and $\mathbb{Z}[B]$-module,
so $H_n(A; \mathbb{Z}[Q]) = H_n(B; \mathbb{Z}[Q]) = 0$ for $n \ge 1$.
The Mayer--Vietoris sequence for $W = A * B$ with coefficients
$\mathbb{Z}[Q]$ collapses to
\[
0 \to H_1(W; \mathbb{Z}[Q])
   \to \mathbb{Z}[Q]/I_A \oplus \mathbb{Z}[Q]/I_B
   \to \mathbb{Z}[Q]/I_W \to 0,
\]
and Crowell's exact sequence identifies
$T_\rho^{\mathrm{ab}} \cong H_1(W; \mathbb{Z}[Q]) \cong I_A \cap I_B$
as a left $\mathbb{Z}[Q]$-submodule.

\emph{Membership.}  $c = XX' \in A$ gives $c - 1 \in I_A$; centrality
of $c$ yields $h = (c-1)(Y-1) = (Y-1)(c-1) \in \mathbb{Z}[Q]\cdot(c-1)
\subseteq I_A$.  Likewise $Y - 1 \in I_B$ and
$h \in \mathbb{Z}[Q]\cdot(Y-1) \subseteq I_B$, so $h \in I_A \cap I_B$.

\emph{Annihilator.}  $(Y+1)(Y-1) = Y^2 - 1 = 0$ and $c$ central, so
$(Y+1) h = (c-1)(Y+1)(Y-1) = 0$.  Similarly $(c+1)(c-1) = 0$ gives
$(c+1) h = 0$.

\emph{Cyclic identification.}  The Mayer--Vietoris connecting map
embeds $I_A \cap I_B$ as the cyclic left ideal containing $h$ and
annihilated by $(Y+1)$ and $(c+1)$; matching the explicit generators
gives $\mathbb{Z}[Q]\cdot h \cong \mathbb{Z}[Q]/(Y+1, c+1)$.
\qed
\end{proof}

\begin{remark}[Centrality is essential]
\label{rem:univ-centrality}
The cyclic identification works precisely because $c \in Q$ is central.
The naive candidate $(X-1)(Y-1)$ fails because $X$ and $Y$ do not
commute in $D_\infty$.  The central element $c$ supplied by the second
$\mathbb{Z}/2$ factor of $Q$ unlocks the cyclic presentation, and is
intrinsic to the right-triangle structure: the central involution
$c = X X'$ realises the half-turn $R_0 R_1$ around the right-angle
vertex.
\end{remark}

\begin{lemma}[Dilation invariance]
\label{lem:dilation-invariance}
For $\lambda > 0$, $\rho_{\lambda T}$ is conjugate to $\rho_T$ by the
central dilation $D_\lambda$, which acts trivially on the
$\mathbb{Z}[Q]$-module structure of $M = \mathbb{Z}[Q]/(Y+1, c+1)$.
\end{lemma}

\begin{proof}
$D_\lambda \circ \rho_T(R_i) \circ D_\lambda^{-1} = \rho_{\lambda T}(R_i)$
for $i = 0, 1, 2$, and $D_\lambda$ has trivial linear part in $Q$.
The defining relations $Y v = -v$, $c v = -v$ of $M$ are
$\mathbb{R}^\times$-equivariant, so the module structure descends.
\qed
\end{proof}

\begin{lemma}[Abstract realization]
\label{lem:rho-abs}
Define $\rho_{\mathrm{abs}} : W \to Q \ltimes M$ by
$R_0 \mapsto (X, 0)$, $R_1 \mapsto (X', 0)$, $R_2 \mapsto (Y, h)$.
Then $\rho_{\mathrm{abs}}$ is a well-defined group homomorphism.
\end{lemma}

\begin{proof}
$\rho_{\mathrm{abs}}(R_0)^2 = (X^2, 0) = (1, 0)$ and similarly for $R_1$;
$\rho_{\mathrm{abs}}(R_2)^2 = (Y^2, (Y+1) h) = (1, 0)$ by Lemma~\ref{lem:trans-mod};
$((X,0)(X',0))^2 = (c^2, 0) = (1, 0)$.
\qed
\end{proof}

\begin{proof}[Proof of Theorem~\ref{thm:universality}]
By Lemma~\ref{lem:trans-mod}, $T_{\rho_T}^{\mathrm{ab}} \cong M$
canonically and independently of $T$; by Lemma~\ref{lem:dilation-invariance}
the isomorphism descends through dilation; by Lemma~\ref{lem:rho-abs}
$\rho_{\mathrm{abs}}$ is well-defined, and every $\rho_T$ factors as
$\rho_T = (\mathrm{id}_Q \ltimes \Phi_T) \circ \rho_{\mathrm{abs}}$ with
$\Phi_T$ a bijection of underlying sets.  Hence $\ker(\rho_T) = \ker(\rho_{\mathrm{abs}})$
for every $T$, and $u_d^T$ is constant in $T$.
\qed
\end{proof}

\section{Asymptotic ratio (table of values)}
\label{sec:tail}

Exact-rational BFS through depth $42$ on the reference $(3, 4, 5)$
triangle yields the values in Table~\ref{tab:tail}.  The asymptotic
behaviour of $u_d$ and the algebraic nature of its generating function
are the principal open questions of the program; partial-progress
results on conditional transcendence, $p$-DFAO state-space analysis,
and structural recurrence exclusions appear in the companion document
\texttt{paper\_extra.tex}.

\begin{table}[ht!]
\centering
\small
\begin{tabular}{r r c r r c}
\toprule
$d$ & $u_d$ & ratio & $d$ & $u_d$ & ratio \\
\midrule
0  & 1                  & ---   & 22 & 41\,067          & 1.515 \\
1  & 3                  & ---   & 23 & 62\,263          & 1.516 \\
2  & 5                  & 1.667 & 24 & 94\,622          & 1.520 \\
3  & 8                  & 1.600 & 25 & 143\,881         & 1.521 \\
4  & 13                 & 1.625 & 26 & 217\,101         & 1.509 \\
5  & 21                 & 1.615 & 27 & 327\,832         & 1.510 \\
6  & 34                 & 1.619 & 28 & 495\,443         & 1.511 \\
7  & 55                 & 1.618 & 29 & 749\,195         & 1.512 \\
8  & 89                 & 1.618 & 30 & 1\,127\,236      & 1.505 \\
9  & 144                & 1.618 & 31 & 1\,697\,179      & 1.506 \\
10 & \textbf{225}       & 1.563 & 32 & 2\,554\,961      & 1.505 \\
11 & 351                & 1.560 & 33 & 3\,848\,384      & 1.506 \\
12 & 554                & 1.578 & 34 & 5\,777\,651      & 1.501 \\
13 & 875                & 1.580 & 35 & 8\,679\,441      & 1.502 \\
14 & 1345               & 1.537 & 36 & 13\,031\,206     & 1.502 \\
15 & 2066               & 1.536 & 37 & 19\,574\,659    & 1.502 \\
16 & 3203               & 1.550 & 38 & 29\,338\,781    & 1.499 \\
17 & 4971               & 1.552 & 39 & 43\,997\,388    & 1.500 \\
18 & 7574               & 1.524 & 40 & 65\,932\,461    & 1.499 \\
19 & 11\,543            & 1.524 & 41 & 98\,849\,591    & 1.499 \\
20 & 17\,683            & 1.532 & 42 & 147\,969\,934   & 1.497 \\
21 & 27\,108            & 1.533 &    &                  &       \\
\bottomrule
\end{tabular}
\caption{The universal sequence $u_d$ through depth $42$.}
\label{tab:tail}
\end{table}

\section{$n$-dimensional extension: upper bound, closed-form rate, and unconditional Class~C faithfulness}
\label{sec:ortho}

The right-triangle analogue in $\mathbb{R}^n$ is the \emph{right-corner
orthoscheme}: a simplex with $n$ mutually perpendicular edges meeting
at a single vertex.  The reflection group $W_n$ is generated by the
$n+1$ face reflections.

\subsection{The Coxeter group is independent of legs}

\begin{lemma}[Perpendicular face pairs, verified symbolically]
\label{lem:perp-pairs}
For the $n$-dim right-corner orthoscheme with vertices on the staircase
$V_0 = 0$, $V_i = V_{i-1} + a_i e_i$, the perpendicular face pairs are
exactly $\{(i, j) : |i - j| \ge 2\}$, independent of the positive leg
values $(a_1, \ldots, a_n)$.  In 3D the dihedral cosines are
$\cos\angle(F_0, F_1) = b/\sqrt{a^2 + b^2}$,
$\cos\angle(F_1, F_2) = ac/\sqrt{(a^2+b^2)(b^2+c^2)}$,
$\cos\angle(F_2, F_3) = b/\sqrt{b^2+c^2}$.
\end{lemma}

The abstract Coxeter group is therefore the right-angled Coxeter group
with edges $(i, j)$ for $|i - j| \ge 2$ and $m_{ij} = \infty$ for
$|i - j| = 1$, the \emph{same group} for every orthoscheme.

\subsection{Steinberg upper bound and closed-form $r_n$}

The commutation graph $G_n$ of the $n$-dim orthoscheme has vertex set
$\{0, \ldots, n\}$ and edges $(i, j)$ when $|i - j| \ge 2$, with
$k$-clique count $\binom{n-k+2}{k}$.  In 3D, Steinberg's formula gives
\[
W(t) \;=\; \frac{(1+t)^2}{1 - 2t},
\]
with $[t^d] W(t) = 9 \cdot 2^{d-2}$ for $d \ge 2$, $u_0 = 1$, $u_1 = 4$,
$u_2 = 9$.  Since $u_d \le [t^d] W(t)$ unconditionally:

\begin{theorem}[Unconditional 3D upper bound]
\label{thm:upper-3d}
For every right-corner orthoscheme in $\mathbb{R}^3$ with positive legs
$(a, b, c)$, $u_d \le 9 \cdot 2^{d-2}$ for $d \ge 2$.
\end{theorem}

\paragraph{The Coxeter diagram.}
Before stating the closed form, we identify the Coxeter diagram
explicitly.  The right-corner orthoscheme reflection group $W_n$ in
$\mathbb{R}^n$ with mutually perpendicular legs $a_1, \ldots, a_n$ has
Coxeter diagram a \emph{path} on $n+1$ nodes (one node per face
reflection).  The leftmost edge --- between the two reflections through
the faces opposite the two legs of greatest aspect ratio --- carries
label $\infty$ (the corresponding dihedral angle is an irrational
multiple of $\pi$ by Lemma~\ref{lem:perp-pairs} together with
Niven--Conway--Jones); the remaining $n-1$ edges carry the standard
label $3$ (a single non-perpendicular pair adjacent to a chain of
right-angle pairs).  This is precisely the path diagram $A_{n+1}^\infty$
for which Cannon--Floyd--Parry, Cannon--Wagreich, and Floyd--Plotnick
compute the growth rate $1 + 2\cos(\pi/(h+1))$ with $h$ the Coxeter
number of the underlying $A_{n+1}$.  Substituting $h = n+2$ gives the
formula of Theorem~\ref{thm:ndim-rate} below.

\begin{theorem}[Closed-form $n$-dim Coxeter rate]
\label{thm:ndim-rate}
For every $n \ge 2$, the abstract orbit-growth rate of the $n$-dim
right-corner orthoscheme reflection group is
\[
r_n  \;=\;  1 \,+\, 2 \cos\!\left(\frac{2\pi}{n+3}\right).
\]
\end{theorem}

\begin{proof}
By Steinberg's formula applied to $G_n$,
$1 / W_n(t^{-1}) = \sum_k (-1)^k \binom{n-k+2}{k} (1+t)^{-k}$.
Substituting $y = -1/(1+t)$ collapses the sum to the Fibonacci--Pell
polynomial $P_{n+3}(y) = \sum_k \binom{n-k+2}{k} y^k$, defined by
$P_0 = 0$, $P_1 = 1$, $P_m = P_{m-1} + y P_{m-2}$.  The characteristic
equation $x^2 = x + y$ has roots whose ratio $(a/b)^m = 1$ gives
$y_j = -1/(4 \cos^2(\pi j / m))$ for $j = 1, \ldots, \lfloor(m-1)/2\rfloor$.
Translating back via $y = -1/(1+t)$ and using
$4\cos^2\theta - 1 = 1 + 2\cos(2\theta)$:
\[
t_j  \;=\;  1 + 2 \cos\!\left(\frac{2\pi j}{n+3}\right);
\]
the asymptotic rate is the maximum, attained at $j = 1$.
\qed
\end{proof}

\begin{corollary}
$3 - r_n = 4\sin^2(\pi/(n+3)) \sim 4\pi^2/(n+3)^2$, so $r_n \to 3$.
\end{corollary}

\begin{remark}[Dynamical-degrees connection]
\label{rem:mcmullen-salem-precise}
The family $\{1 + 2\cos(2\pi/m)\}_{m \ge 3}$ also arises in birational
surface dynamics (McMullen~2007; Bedford--Kim 2009).  For $n \ne 3$,
$r_n$ is Pisot or totally-real of degree $\ge 3$, hence \emph{not} a
Salem number, and by McMullen~2007 Theorem~1.5 cannot be the dynamical
degree of a projective rational-surface automorphism.  The case
$r_3 = 2$ is realised as McMullen's $\sigma_9$.
\end{remark}

\subsection{Class C unconditional faithfulness}

\begin{lemma}[Uniform Schur-complement determinant identity]
\label{lem:uniform-detQ}
For every $n \ge 3$ and every positive tuple $(a_1, \ldots, a_n)$, the
$(n+1) \times (n+1)$ Lorentzian Gram matrix $Q_n$ of the $n$-dim
right-corner orthoscheme satisfies
\[
\det Q_n \;=\; -\, \prod_{i=1}^{n} a_i^2.
\]
Here $Q_n$ is the symmetric tridiagonal matrix with
$Q_n[0,0] = 1 - a_1^2$, $Q_n[i,i] = a_i^2 + a_{i+1}^2$ ($1 \le i \le n-1$),
$Q_n[n,n] = 1$, $Q_n[0,1] = a_2$, $Q_n[i,i+1] = a_i a_{i+2}$ ($1 \le i \le n-2$),
$Q_n[n-1,n] = a_{n-1}$, read from the staircase normal vectors with the
Lorentzian endpoint correction.
\end{lemma}

\begin{proof}
Let $D_k = \det Q_n[:k{+}1,:k{+}1]$.  Tridiagonal expansion gives
$D_k = Q_n[k,k] D_{k-1} - Q_n[k-1,k]^2 D_{k-2}$ with $D_{-1} = 1$.
We claim $D_k = P_k := (\prod_{i=1}^k a_i^2)(1 - \sum_{i=1}^{k+1} a_i^2)$
for $1 \le k \le n-1$, and $D_n = -\prod_{i=1}^n a_i^2$.

\emph{Base.}  $D_0 = 1 - a_1^2$; $D_1 = (1-a_1^2)(a_1^2+a_2^2) - a_2^2
= a_1^2(1 - a_1^2 - a_2^2) = P_1$.

\emph{Interior step.}  For $2 \le k \le n-1$, the recurrence reduces to
the polynomial identity (in the symbols $a_{k-1}, a_k, a_{k+1}, S_{k-1}, P_{k-2}$
with $S_{k-1} = \sum_{i=1}^{k-1} a_i^2$, $P_{k-2} = \prod_{i=1}^{k-2} a_i^2$)
\[
P_{k-2}\, a_{k-1}^2 a_k^2\, (1 - S_{k-1} - a_k^2 - a_{k+1}^2)
=
(a_k^2 + a_{k+1}^2)\, P_{k-2}\, a_{k-1}^2\, (1 - S_{k-1} - a_k^2)
- (a_{k-1} a_{k+1})^2\, P_{k-2}\, (1 - S_{k-1}),
\]
which expands to zero in $\mathbb{Z}[a_{k-1}, a_k, a_{k+1}, S_{k-1}, P_{k-2}]$
and is verified symbolically.  The identity is independent of $k$ and
of $n$.  The base, inductive, and boundary steps are formally verified
in Lean~4 (theorems \texttt{schur\_base\_step}, \texttt{schur\_k2\_step},
\texttt{schur\_inductive\_step}, \texttt{schur\_boundary\_step} in
\texttt{lean/with\_mathlib/SchurGeneral.lean}, each discharged by the
\texttt{ring} tactic), and the determinant identity
$\det Q_n = -\prod_{i=1}^n a_i^2$ itself is machine-checked for
$n = 2, \dots, 10$ on explicit leg sequences in
\texttt{lean/RightTriangleReflection.lean}.

\emph{Final step.}  At $k = n$, $D_n = D_{n-1} - a_{n-1}^2 D_{n-2}$, and
\[
D_{n-1} - a_{n-1}^2 D_{n-2}
= P_{n-2}\, a_{n-1}^2 \bigl[(1 - S_{n-1} - a_n^2) - (1 - S_{n-1})\bigr]
= -\prod_{i=1}^n a_i^2.
\]
The three identities (base, interior, final) close the induction for
every $n \ge 3$.
\qed
\end{proof}

\begin{theorem}[Class C unconditional faithfulness, all $n \ge 3$]
\label{thm:master-C-uncond}
For every $n \ge 3$ and every leg tuple $(a_1, \ldots, a_n) \in \mathbb{Q}_{>0}^n$
with pairwise distinct coordinates, the affine representation
$\rho : W_n \to \mathrm{Aff}(\mathbb{R}^n)$ of the $n$-dim right-corner
orthoscheme reflection group is faithful, and
\[
u_d(O) \;=\; [t^d]\, W_n(t)
\qquad \text{for every } d \ge 0.
\]
In particular for $n = 3$ this is $u_d = 9 \cdot 2^{d-2}$ for $d \ge 2$.
\end{theorem}

\begin{proof}
Write $S = \{R_0, \ldots, R_n\}$ for the face reflections and let
$m_{ij} \in \{2, 3, \ldots, \infty\}$ be the order of $R_i R_j$ in the
group $W_n^{\mathrm{geo}} = \langle S\rangle$ actually generated inside
$\mathrm{Isom}(\mathbb{R}^n)$.  Faithfulness of $\rho$ together with the
identity $u_d(O) = [t^d]\, W_n(t)$ is equivalent to the assertion that
$(m_{ij})$ equals the intended right-angled matrix
\begin{equation}
m_{ij} = 2 \ (|i-j| \ge 2), \qquad m_{i,i+1} = \infty \ (0 \le i \le n-1),
\tag{$\ast$}\label{eq:intended-coxeter}
\end{equation}
i.e.\ that the orthoscheme introduces \emph{no} relation $(R_i R_j)^m = 1$
beyond the perpendicular pairs.  We prove~\eqref{eq:intended-coxeter}; the
theorem then follows from Tits' theorem applied to the resulting Coxeter
system (end of proof).  We emphasize that non-degeneracy of $Q_n$
(Lemma~\ref{lem:uniform-detQ}) does \emph{not} by itself give
faithfulness: a non-degenerate Gram form is compatible with a rational
dihedral angle.  Indeed the isosceles case $a = b$ has
$\det Q_n = -\prod a_i^2 \ne 0$ yet $\theta_{01} = \pi/4$, giving the
extra relation $(R_0 R_1)^4 = 1$ --- the $\mathrm{cd} = 4$ case of
Theorem~\ref{thm:cd-general}.  The content of the proof is therefore the
angle-irrationality argument that follows; $\det Q_n \ne 0$ enters only
to confirm the simplex is non-degenerate for all positive legs.

\emph{Step 1 (dihedral angles $\leftrightarrow$ relations).}  For two
faces meeting along a codimension-$2$ ridge, $R_i, R_j$ generate a finite
dihedral group of order $2m$ iff the dihedral angle $\theta_{ij} = \pi/m$
for an integer $m \ge 2$, and an infinite dihedral group (no relation)
iff $\theta_{ij} \notin \pi\mathbb{Q}$.  By Lemma~\ref{lem:perp-pairs}
the pairs with $|i-j| \ge 2$ are perpendicular ($m_{ij} = 2$), giving the
right-angled edges of~\eqref{eq:intended-coxeter}.  It remains to show
every non-perpendicular $\theta_{i,i+1}$ is an irrational multiple
of~$\pi$.

\emph{Step 2 (Niven--Conway--Jones constraint).}  By
Lemma~\ref{lem:perp-pairs} each non-perpendicular $\cos^2\theta_{i,i+1}$
is rational in the legs.  If some $\theta_{i,i+1} \in \pi\mathbb{Q}$,
then by the extended Niven theorem (Conway--Jones 1976) a rational angle
with rational squared cosine satisfies $\cos\theta_{i,i+1} \in
\{0, \pm\tfrac12, \pm\tfrac{\sqrt2}{2}, \pm\tfrac{\sqrt3}{2}, \pm 1\}$,
i.e.\ $\cos^2\theta_{i,i+1} \in \{0, \tfrac14, \tfrac12, \tfrac34, 1\}$.
The values $0$ (perpendicular) and $1$ (degenerate simplex) are
excluded, so a genuine extra relation requires
$\cos^2\theta_{i,i+1} \in \{\tfrac14, \tfrac12, \tfrac34\}$, i.e.\
$\theta_{i,i+1} \in \{\pi/3, \pi/4, \pi/6\}$.

\emph{Step 3 (boundary dihedrals; elementary).}  The end edges have
$\cos^2\theta_{01} = b^2/(a^2+b^2)$ in the local two-leg notation
$(a, b) = (a_1, a_2)$, and symmetrically at the far end.  Setting this
equal to $\tfrac14, \tfrac12, \tfrac34$ forces $b^2 = \tfrac13 a^2$,
$b^2 = a^2$, $b^2 = 3 a^2$ respectively: the first and third give
$b/a \in \{1/\sqrt3, \sqrt3\} \notin \mathbb{Q}$, contradicting
rationality of the legs; the middle gives $a = b$, excluded by the
pairwise-distinct hypothesis.  No boundary dihedral is a rational
multiple of $\pi$, and no elliptic curve is needed here.

\emph{Step 4 (interior dihedrals; square-discriminant $\to$ Mordell).}
An interior edge involves three consecutive legs $(a, b, c) =
(a_i, a_{i+1}, a_{i+2})$ with $b$ the middle leg, and by
Lemma~\ref{lem:perp-pairs}
\[
\cos^2\theta_{i,i+1} = \frac{a^2 c^2}{(a^2+b^2)(b^2+c^2)}.
\]
Setting this equal to a Niven value $\tfrac14, \tfrac12, \tfrac34$ and
clearing denominators yields a quadratic in $B := b^2$.  Since the legs
are rational, $B = b^2 \in \mathbb{Q}$, so this quadratic has a rational
root only if its discriminant is a rational square.  The three
discriminants are, up to a positive rational square factor, exactly the
Diophantine varieties
\[
\mathcal{V}_{1/4}:\ a^4 + 14 a^2 c^2 + c^4 = m^2,
\quad
\mathcal{V}_{1/2}:\ a^4 + 6 a^2 c^2 + c^4 = k^2,
\quad
\mathcal{V}_{3/4}:\ 9 a^4 + 30 a^2 c^2 + 9 c^4 = k^2
\]
in the two flanking legs (this square-discriminant condition is the
correct origin of the three varieties).  A rational interior angle thus
forces a rational point on one of $\mathcal{V}_{1/4}, \mathcal{V}_{1/2},
\mathcal{V}_{3/4}$ with $a, c \in \mathbb{Q}_{>0}$.  Each is a genus-$1$
curve in $t = a/c$ whose Jacobian has rank $0$, so each has only finitely
many rational points:
\begin{itemize}[leftmargin=*,itemsep=0.1em]
\item $\mathcal{V}_{1/4}$: Jacobian $y^2 = x(x-12)(x-16)$, Cremona
\texttt{24a1} (conductor $24$, rank $0$, torsion
$\mathbb{Z}/4 \oplus \mathbb{Z}/2$); rational points $t = a/c \in
\{-1, 0, 1\}$.
\item $\mathcal{V}_{1/2}$: Jacobian $y^2 = x^3 - x$, Cremona
\texttt{32a2} (conductor $32$, rank $0$, torsion
$\mathbb{Z}/2 \oplus \mathbb{Z}/2$, the four points
$(0,0), (\pm 1, 0), \infty$); rational points $t = 0$.
\item $\mathcal{V}_{3/4}$: Jacobian $y^2 = x(x-300)(x-1200)$, Cremona
\texttt{72a2} (conductor $72$, rank $0$, torsion
$\mathbb{Z}/2 \oplus \mathbb{Z}/2$); rational points $t = 0$.
\end{itemize}
Every such point gives $a c = 0$ or $a = \pm c$, all excluded by the
pairwise-distinct, strictly positive hypothesis on the legs.  (Ranks,
torsion, Cremona labels, and the rational-point determinations are
verified in SageMath in \texttt{code/mordell/mordell\_24a1.sage},
\texttt{mordell\_32a2.sage}, \texttt{mordell\_72a2.sage}, with output
committed in the project repository.)  Hence no interior dihedral is a
rational multiple of $\pi$ either.

\emph{Conclusion.}  By Steps 3--4 no non-perpendicular $\theta_{i,i+1}$
lies in $\pi\mathbb{Q}$, so $m_{i,i+1} = \infty$, and with Step 1 the
geometric Coxeter matrix is exactly~\eqref{eq:intended-coxeter}.  Thus
$W_n^{\mathrm{geo}}$ is the Coxeter system with matrix
\eqref{eq:intended-coxeter}, and Tits' theorem (Humphreys,
\emph{Reflection Groups and Coxeter Groups}, Thm.~5.4) applied to
\emph{this established system} gives faithfulness of the geometric
representation $\rho$.  Consequently
$u_d(O) = \#\{w : \ell(w) = d\} = [t^d]\, W_n(t)$; for $n = 3$ this is
$9 \cdot 2^{d-2}$ for $d \ge 2$.
\qed
\end{proof}

\subsection{The collision-depth invariant $\mathrm{cd}_n \in \{\infty, 4, 3\}$}

For a leg sequence $(a_1, \ldots, a_n)$ define
\[
e \;:=\; \#\bigl\{i \in \{1,\, n-1\} : a_i = a_{i+1}\bigr\} \in \{0, 1, 2\}, \qquad
t \;:=\; \#\bigl\{i : 1 \le i \le n-2,\ a_i = a_{i+1} = a_{i+2}\bigr\}.
\]
$e$ counts \emph{endpoint pairs} (adjacent equal legs at the leftmost
or rightmost positions) and $t$ counts \emph{triples}.  Let
$\mathrm{cd}_n(a_1, \ldots, a_n)$ be the minimal word-length at which
the orbit BFS first exhibits a non-trivial collision (i.e.\ the first
depth at which $u_d < [t^d] W_n(t)$); set $\mathrm{cd}_n = \infty$ when
the affine representation is faithful at every length.

\begin{theorem}[Collision-depth invariant, all $n \ge 3$]
\label{thm:cd-general}
For every $n \ge 3$ and every $(a_1, \ldots, a_n) \in \mathbb{Q}_{>0}^n$,
\[
\mathrm{cd}_n(a_1, \ldots, a_n) \;=\;
\begin{cases}
\infty & \text{if } e = t = 0, \\
4 & \text{if } e \ge 1,\ t = 0, \\
3 & \text{if } t \ge 1.
\end{cases}
\]
\end{theorem}

\begin{proof}
Faithfulness for $e = t = 0$ is Theorem~\ref{thm:master-C-uncond}.  An
endpoint pair $a_i = a_{i+1}$ at a boundary makes the corresponding
boundary cosine $b/\sqrt{a^2+b^2}$ equal $\pm\sqrt 2/2$, producing the
length-$4$ relation $(R_i R_{i+1})^4 = \mathrm{id}$ and hence
$\mathrm{cd} \le 4$; no shorter relation can appear because
$\cos^2\theta = 0, 1$ are degenerate and the $\pm 1/2$, $\pm\sqrt 3/2$
cases reduce as in the proof of Theorem~\ref{thm:master-C-uncond} to
empty or degenerate loci.  An interior triple $a_i = a_{i+1} = a_{i+2}$
makes the middle cosine $ac/\sqrt{(a^2+b^2)(b^2+c^2)} = 1/2$ (after the
quadratic-in-$b^2$ analysis with $a = b = c$), producing the length-$3$
braid relation $R_i R_{i+1} R_i = R_{i+1} R_i R_{i+1}$ and hence
$\mathrm{cd} \le 3$.  An interior pair (with no endpoint pair and no
triple) leaves all flanking dihedrals irrational multiples of $\pi$,
since the Niven exceptional values reduce to the same three Diophantine
varieties $\mathcal{V}_{1/4}, \mathcal{V}_{1/2}, \mathcal{V}_{3/4}$
which are empty over distinct-coordinate $\mathbb{Q}_{>0}^2$ (proof of
Theorem~\ref{thm:master-C-uncond}); hence $\mathrm{cd} = \infty$.
\qed
\end{proof}

Note that in $n = 3$ every adjacent equal pair is automatically an
endpoint pair, so the formula reduces to the run-length statistic; in
$n \ge 4$ \emph{interior} pairs $a_i = a_{i+1}$ (e.g.\ $a_2 = a_3$ in
$(1, 2, 2, 3)$) leave $\mathrm{cd} = \infty$, distinguishing $(e, t)$
sharply from the longest-run statistic.

\section{Reproducibility}
\label{sec:reproduce}

All computations use exact rational arithmetic; the symbolic
verification of Theorem~\ref{thm:symbolic} uses SymPy in
$\mathbb{Q}(a, b)$, and the Schur-complement identity of
Lemma~\ref{lem:uniform-detQ} is verified symbolically in
$\mathbb{Z}[a_{k-1}, a_k, a_{k+1}, S_{k-1}, P_{k-2}]$.  The rank-$0$
descents on the Cremona curves \texttt{24a1}, \texttt{32a2}, and
\texttt{72a2} (for $\mathcal{V}_{1/4}, \mathcal{V}_{1/2},
\mathcal{V}_{3/4}$ respectively) are verified in SageMath via
\texttt{EllipticCurve.rank()} and a direct rational-point determination,
in \texttt{code/mordell/}.  The BFS values in
Table~\ref{tab:tail} are computed by an exact-rational BFS through
depth $25$ and a Rust disk-streaming BFS through depth $42$.  Source
code and reproduce scripts accompany this paper in the project
repository.

\paragraph{Formal (machine-checked) verification.}
Several of the computational claims are additionally certified in the
Lean~4 proof assistant, in the \texttt{lean/} subtree of the project
repository.  A core, Mathlib-free development
(\texttt{RightTriangleReflection.lean}) machine-checks the basic Coxeter
relations, all eight length-$10$ relations and their explicit affine
matrices on the $(3,4,5)$ triangle, a direct breadth-first reconstruction
of $a(0),\dots,a(17)$, the Fibonacci coincidence $a(n) = F_{n+3}$ for
$1 \le n \le 9$, and the determinant identity
$\det Q_n = -\prod a_i^2$ for $n = 2,\dots,10$.  A second development
depending on Mathlib (\texttt{with\_mathlib/}) proves the eight length-$10$
relations symbolically over $\mathbb{Q}(a,b)$
(Theorem~\ref{thm:symbolic}), the algebraic steps of the Schur induction
(Lemma~\ref{lem:uniform-detQ}), and the universality of the orbit-growth
sequence through depth $22$ (Remark~\ref{rem:lean-universality}), the
last holding simultaneously for all right triangles with unequal legs.
The \texttt{ring}-based theorems (the symbolic length-$10$ relations and
the Schur induction steps) are verified entirely by the Lean kernel; the
remaining results use \texttt{native\_decide}, whose trusted base
additionally includes the Lean compiler and the correctness of the
hand-rolled \texttt{partial} functions for polynomial determinant and
deduplication.  Building the projects re-verifies every claim.

\section*{References}

\begin{itemize}[leftmargin=2em,itemsep=0.2em]
\item J.~E.~Humphreys, \emph{Reflection Groups and Coxeter Groups},
Cambridge Studies in Advanced Mathematics 29, Cambridge University
Press, 1990.  Tits' faithfulness theorem (Thm.~5.4) and Steinberg's
word-growth formula (\S5.12).

\item I.~M.~Niven, \emph{Irrational Numbers}, Carus Mathematical
Monographs 11, MAA, 1956.

\item J.~H.~Conway, A.~J.~Jones, \emph{Trigonometric Diophantine
equations (On vanishing sums of roots of unity)}, Acta Arith.\ 30
(1976), 229--240.

\item J.~Steinberg, \emph{Endomorphisms of Linear Algebraic Groups},
Memoirs of the AMS 80, 1968.

\item L.~J.~Mordell, \emph{Diophantine Equations}, Pure and Applied
Mathematics 30, Academic Press, 1969.  §17 (rank-$0$ descent on
$x^4 + 6 x^2 y^2 + y^4 = z^2$).

\item J.~W.~S.~Cassels, \emph{Lectures on Elliptic Curves}, LMS Student
Texts 24, Cambridge University Press, 1991.  Ch.~II on Fermat-style
descent.

\item J.~E.~Cremona, \emph{Algorithms for Modular Elliptic Curves},
2nd edition, Cambridge University Press, 1997; online tables at
\url{https://www.lmfdb.org/EllipticCurve/Q/}.  Curves
\href{https://www.lmfdb.org/EllipticCurve/Q/24/a/1}{\texttt{24a1}} and
\href{https://www.lmfdb.org/EllipticCurve/Q/72/a/2}{\texttt{72a2}}.

\item M.~W.~Davis, \emph{The Geometry and Topology of Coxeter Groups},
LMS Monographs 32, Princeton University Press, 2008.

\item J.~W.~Cannon, P.~Wagreich, \emph{Growth functions of surface
groups}, Math.\ Ann.\ 293 (1992), no.~2, 239--257.

\item W.~J.~Floyd, S.~P.~Plotnick, \emph{Growth functions on Fuchsian
groups and the Euler characteristic}, Invent.\ Math.\ 88 (1987),
no.~1, 1--29.

\item H.~S.~M.~Coxeter, \emph{Discrete groups generated by
reflections}, Annals of Math.\ 35 (1934), 588--621.

\item C.~T.~McMullen, \emph{Dynamics on blowups of the projective
plane}, Publ.\ Math.\ IHES 105 (2007), 49--89.

\item E.~Bedford, K.~Kim, \emph{Dynamics of rational surface
automorphisms: linear fractional recurrences}, J.\ Geom.\ Anal.\ 19
(2009), no.~3, 553--583.

\item J.~Diller, C.~Favre, \emph{Dynamics of bimeromorphic maps of
surfaces}, Amer.\ J.\ Math.\ 123 (2001), no.~6, 1135--1169.

\item OEIS Foundation, \emph{The On-Line Encyclopedia of Integer
Sequences}, \url{https://oeis.org}.  Sequence A000045 (Fibonacci
numbers).
\end{itemize}

\end{document}
